An Arrow-Debreu model concerns itself with the interactions of
decentralized decisions by consumers and producers in markets.
\citet{Lars} proposed a representation of this class of models in
which two types of equations define an equilibrium: zero profit and
market clearance. The corresponding variables defining an equilibrium
are activity levels (for constant-returns-to-scale firms) and
commodity prices.\footnote{Under a maintained assumption of perfect
competition, \citet{Lars} may characterize technology as
constant-returns-to-scale without loss of generality. Specifically,
decreasing returns are accommodated through introduction of a specific
factor, while increasing returns are inconsistent with the assumption
of perfect competition. Note that in this environment zero excess
profit is consistent with free entry for atomistic firms producing an
identical product.} Here we extend Mathiesen's framework with a third
class of variables corresponding to consumer income levels. Commodity
markets encompass primary endowments of households with producer
outputs. In equilibrium the aggregate supply of each good must be at
least as great as total intermediate and final demand. Initial
endowments are exogenous. Producer supplies and demands are defined by
producer activity levels and relative prices. Final demands are
determined by market prices.

Economists who have worked with conventional textbook equilibrium
models can find Mathiesen's framework to be somewhat opaque because
many quantity variables need not be explicitly specified in the model.
Variables such as final demand by consumers, factor demands by
producers and commodity supplies by producers, are defined implicitly
in Mathiesen's model. For example, given equilibrium prices for
primary factors, consumer incomes can be computed, and given income
and goods prices, consumers' demands can then be determined. The
consumer demand functions are written down in order to define an
equilibrium, but quantities demanded need not appear in the model as
separate variables. The same is true of inputs or outputs from the
production process: relative prices determine conditional demand, and
conditional demand times the activity level represents market demand.
Omitting decisions variables and suppressing definitional equations
corresponding to intermediate and final demand provides significant
computational advantages at the cost of a somewhat more complicated
model statement.

In the following, we detail (i) zero profit conditions, (ii) market
clearance conditions, and (iii) income balance conditions, which in
the present case is equivalent to the regional budget constraint.
These three sets of conditions form the basic system of equation to be
solved. Note that the actual code for the model is implemented both in
the algebraic mixed complementarity format \citep[\GAMS/\MCP,
see][]{Rutherford1995} and through the more compact formulation
afforded by the Mathematical Programming System for General
Equilibruim (\MPSGE) syntax \citep{Rutherford1999}.

\subsection[Zero profit]{Zero profit (arbitrage) conditions}

All production activities in the model are represented by
constant-returns-to-scale technologies, and markets are assumed to
operate competitively with free entry and exit. As a consequence,
equilibrium profits are driven to zero and the price of output
reflects the cost of inputs. The following sets of equations relating
output price to marginal cost are part of the definition of an
equilibrium.\footnote{To retain consistency with the \MCP format, we
express zero profit conditions as ``oriented equations,'' with
marginal (=average) cost on the \LHS and marginal (=average) revenue on
the \RHS.}

The calculation of unit cost and unit revenue functions involves the
definition of a number of ancillary variables (that do not appear in
the \GAMS code as explicit choice variables). In the following we
define ancillary variables in un-numbered equations, indicating that
these variables are ``optional'' in the sense that they may be
substituted out of the non-linear system of equations. Moreover, we
use the symbol $\theta$ to portray value shares from the base year
data. In most cases subscripts on these value shares are omitted in
order to economize on notation. Finally, to denote benchmark values we
use an overline, so that $\overline{t}^{pd}_{ir}$ represents the
benchmark value of ${t}^{pd}_{ir}$.

\subsubsection[Production]{Sectoral production (\texttt{Y(j,r)})}

Sectoral production combines intermediate inputs with a value-added
nest combining primary inputs (see Figure \ref{prodfunY}). The unit
cost of value-added is a \CES composite of skilled and unskilled labor,
land, resources and capital inputs to production, gross of taxes.
Factor inputs may be sector-specific or mobile across sectors:
$$p^{pf}_{fjr} = \left\{\begin{array}{lr}
p^F_{fr} \frac{(1+t^f_{fjr})}{1+\overline{t}^f_{fjr}} & f \in m_f\\
p^S_{fjr} \frac{(1+t^f_{fjr})}{1+\overline{t}^f_{fjr}} & f \in s_f
\end{array} \right.$$
and the unit cost function is given by:
$$c^f_{jr} = \left( \sum_f \theta_f \left( p^{pf}_{fjr}
\right)^{1-\sigma} \right)^{1/(1-\sigma)},$$
in which $\theta_f$ represents the factor share of value added
[{\small \texttt{theta\_vfm(f,g,r)}}].

The user cost of intermediate inputs differs from the market price
due to the presence of taxes on intermediate inputs:
$$p^d_{ijr} = p^Y_{ir} \frac{1+t^{fd}_{ijr}}{1+\overline{t}^{fd}_{ijr}}$$
$$p^i_{ijr} = p^M_{ir}
\frac{1+t^{fi}_{ijr}}{1+\overline{t}^{fi}_{ijr}}$$ A \CES cost function
describes the minimum cost of a bundle of domestic and imported inputs
to production, based on benchmark value shares and an elasticity of
substitution $\sigma [\texttt{esubd(i)}]$:
$$c^i_{ijr} = \left( \theta (p^d_{ijr})^{1-\sigma} +
             (1-\theta) (p^i_{ijr})^{1-\sigma}
\right)^{1/(1-\sigma)},$$
in which $\theta$ represents the domestic share of the Armington
composite  \hfill \break [{\small \texttt{theta\_vdfm(i,j,r)}}].

Unit cost of sectoral output is then a Leontief (linear) composite of
the costs of intermediate and value-added composite inputs, based on
base-year value shares:
$$c^Y_{jr} = \sum_i \theta_i c^i_{ijr} + \theta^f c^f_{jr},$$
in which $\theta_i$ represents the cost share of intermediate input
$i$ [{\small \texttt{theta\_cm0(i,j,r)}}]; and
$\theta^f$ represents the cost share of value added in sectoral
output [{\small \texttt{theta\_cf0(j,r)}}].

Having formulated the unit cost function, it is possible to compactly
portray the zero profit condition for $y_{jr}$. In equilibrium, the
marginal cost of supply equals the market price, net of taxes:
\begin{equation}
c^Y_{jr} = p^Y_{jr} \frac{1-t^o_{jr}}{1-\overline{t}^o_{jr}}
\end{equation}
In the \SOE model production for domestic and export
markets are differentiated, and we replace $p^Y_{jr}$ by a unit
revenue function,
$$r^Y_{jr} = \left( \theta (p^D_{jr})^{1+\eta} +
(1-\theta) (p^X_{jr})^{1+\eta} \right)^{1/(1+\eta)},$$
in which $\theta$ represents the export share of output
[{\small \texttt{theta\_vxm(j,r)}}].

\subsubsection[Private demand]{Consumer demand (\texttt{Y("c",r)})}

In previous versions of \GTAPINGAMS the consumer price index
represented a Cobb-Douglas demand system with an Armington
substitution between domestic and imported goods prices gross of tax.
This model is retained as one option in the new model, price indices
for domestic and imported goods are given by:
$$p^{dc}_{ir} = p^Y_{ir} \frac{1+t^{pd}_{ir}}{1+\overline{t}^{pd}_{ir}},$$
and
$$p^{ic}_{ir} = p^M_{ir} \frac{1+t^{pi}_{ir}}{1+\overline{t}^{pi}_{ir}},$$
and the Armington composite price of good $i$ is a \CES composite price
defined over these indices:
$$p^c_{ir} = \left(\theta (p^{dc}_{ir})^{1-\sigma} +
                (1-\theta) (p^{ic}_{ir})^{1-\sigma}
\right)^{1/(1-\sigma)},$$
in which $\theta$ represents the domestic demand share
[{\small \texttt{theta\_vdfm(i,"c",r)}}].

In addition to the Cobb-Douglas representation of the final demand
system, our core model includes a \LES, which combines a Cobb-Douglas
discretionary demand with a Leontief subsistence demand activity.  The
added complexity of the \LES specification permits the model to be
calibrated both to the income and \textit{average} price elasticities
of demand. (See Appendix \ref{LES} of \texttt{documentation.pdf} for details.) The price indices for
discretionary and subsistence demand are defined by:
\begin{equation}
p^{DD}_r = \prod_i (p^c_{ir})^{\theta_i},
\end{equation}
in which $\theta_i$ represents the discretionary value share
[{\small \texttt{theta\_cm0(i,"dd",r)}}].

The price index for subsistence demand is defined by a Leontief
cost function:
\begin{equation}
p^{SD}_r = \sum_i \theta_i p^c_{ir},
\end{equation}
in which $\theta_i$ represents the commodity $i$ share of subsistence
demand  \hfill \break [{\small \texttt{theta\_cm0(i,"sd",r)}}].

Finally, \GAMS code includes a third representation of final demand in
the form of the \CDE demand system \citep{Hanoch1975}. This
representation, details of which are presented in Appendix \ref{CDE} of \texttt{documentation.pdf},
has been part of the \GTAP model in \textsc{gempack} for many years,
and our implementation follows those of
\citet{HertelvanderMensbrugghe2015} and \citet{Chen}. Similar to the
LES demand system, this approach permits calibrating empirical
evidence on both own-price and income elasticities.

\subsubsection[Government demand]{Government demand (\texttt{Y("g",r)})}

Public expenditure is a fixed-coefficient aggregate of Armington
composite goods. Within each composite domestic and imported goods
trade off with a constant elasticity of substitution. The unit price
indices for domestic and imported goods are given by:
$$p^{dg}_{ir} = p^Y_{ir} \frac{1+t^{dg}_{ir}}{1+\overline{t}^{gd}_{ir}} $$
and
$$p^{ig}_{ir} = p^M_{ir} \frac{1+t^{ig}_{ir}}{1+\overline{t}^{ig}_{ir}}$$
The composite price of the $i$th good is then:
$$pg_{ir} = \left( \theta (p^{dg}_{ir})^{1-\sigma} + (1-\theta)
(p^{ig}_{ir})^{1-\sigma} \right)^{1/(1-\sigma)},$$
in which  $\theta$ represents the domestic share of public demand
[{\small \texttt{theta\_vdfm(i,"g",r)}}].

The unit cost of public services ($G_{r}$) is defined by the Leontief
cost coefficients:
\begin{equation}
\sum_i \theta_i pg_{ir}  = p^G_r,
\end{equation}
in which $\theta_i$ represents the value share of commodity $i$ in
public expenditure \hfill \break [{\small \texttt{theta\_cm0(i,"g",r)}}].

\subsubsection[Aggregate imports]{Aggregate imports (\texttt{M(i,r)})}

An import cost index applies export taxes, trade and transport margins
and import tariffs to the producer supply prices in exporting regions:
$$py^m_{isr} = p^Y_{is}
\frac{(1-t^{xs}_{isr})(1+t^{ms}_{isr})}
{(1-\overline{t}^{xs}_{isr})(1+\overline{t}^{ms}_{isr})}\,.$$ In the
case of the \SOE model, the supply price $p^Y_{is}$ is replaced by
export price $p^X_{is}$. The unit price of transportation services is
given by: $$pt^m_{jisr} = p^T_j
\frac{1+t^{ms}_{isr}}{1+\overline{t}^{ms}_{isr}}\,.$$

Transportation margins enter as fixed coefficients with bilateral
trade flows, so the unit delivered price is a convex combination of
the unit prices with weights corresponding to base year value shares:
$$pyt^m_{isr} = \theta py^m_{isr} + \sum_j \theta^T_j pt^m_{jisr},$$
in which $\theta$ represents good $i$ share of imports
[{\small \texttt{theta\_vxmd(i,s,r)}}]; and $\theta^T_j$
represents the value share of transportation service $j$ in the import
price \hfill \break [{\small \texttt{theta\_vtwr(j,i,s,r)}}].

Having formed a price index for bilateral imports from region $s$ to
region $r$, the \CES cost index can be defined on the basis of value
shares and the elasticity of substitution across imports from
different regions, $\sigma = \textit{esubm}_i$:
$$\textit{cim}_{ir} = \left(\sum_s \theta_s (pyt^m_{isr})^{1-\sigma)}
\right)^{1/(1-\sigma)} ,$$
in which $\theta_s$ is the  value share of bilateral imports from
region $s$ [{\small \texttt{theta\_m(i,s,r)}}].

The import activity ($m_{ir}$) has a zero
profit condition which relates the unit cost of imports to the market
price of the import aggregate:
\begin{equation}
cim_{ir} = p^M_{ir}\,.
\end{equation}

\subsubsection[Transportation]{International transportation services
(\texttt{YT(j)})}

For simplicity, the unit cost of a transportation service depends on
the benchmark value shares of region-specific services through a
Cobb-Douglas cost function. Under perfect competition with free entry,
the unit cost of international transport services equals the
equilibrium market price:
\begin{equation}
\prod_r (p^Y_{jr})^{\theta_r} = p^T_j
\end{equation}
in which $\theta_r$ represents the region $r$ share of transportation
service $j$ [{\small \texttt{theta\_vst(j,r)}}].
In the \SOE model, the supply price $p^Y_{is}$ in equation is replaced
by export price $p^X_{is}$.

\subsubsection[Sector-specific factors]{Sector-specific factor transformation (\texttt{FT(f,r)})}

The unit value of sector-specific factors is defined as a \CET
aggregate of returns to factor $f$ across sectors $j$:
$$\textit{pvfm}_{fr} = \left( \sum_j \theta_j ps_{fjr}^{1+\eta}
\right)^{1/(1+\eta)} \quad f \in s_{f}\,.$$
in which $\theta_j$ is the sector $j$ share of earnings for factor
$f$ [{\small \texttt{theta\_evom(f,j,r)}}].

The constant elasticity of transformation frontier defines the
profit-maximizing allocation of factors to individual sectors. In
equilibrium, the unit value of the aggregate factor is equal to the
maximum unit earnings:
\begin{equation}
p^S_{fr} = \textit{pvfm}_{fr} \quad f \in s_f
\end{equation}

\subsection{Market clearance}

Supply-demand conditions apply to all goods and factors. Benchmark
demand and supply quantities appear as scale factors in many of these
equations, typically multiplied by activity levels which are equal to
unity in the reference equilibrium.\footnote{While not crucial for
representation of the model as a nonlinear system of equations, we
follow the \MCP convention in writing out the market clearance
conditions. The equations are ``oriented'', with supply variables on
the \LHS and demands on the \RHS. Hence, the sense of the equation is
$supply \geq demand$. In the core model equilibrium prices should
always be positive, but in extensions of the standard model it might
be quite common to introduce inequalities and complementary slackness,
in which case the proper orientation of the equations is essential.
Hence, in equilibrium should the price of a good be zero, economic
equilibrium is then consistent with a market in which $supply >
demand$.}

\subsubsection[Firm output]{Firm output (\texttt{P(i,r)})}

Aggregate output of good $i$ in region $r$ in the reference equilibrium is
\texttt{vom(i,r)}:
\begin{equation}
\label{Youtput}
Y_{ir} \textit{vom}_{ir} =
\sum_j \textit{ddfm}_{ijr} + \textit{ddfm}_{iCr} + \textit{ddfm}_{iIr}
+ \textit{ddfm}_{iGr} + \sum_s \textit{dxmd}_{irs} + \textit{dst}_{ir}
\end{equation}
where the compensated demand functions can be obtained by
differentiating the unit cost functions:
\begin{eqnarray}
\textit{ddfm}_{ijr} &=& Y_{jr} \textit{vdfm}_{ijr} \left( \frac{\textit{ci}_{ijr}}{p^d_{ijr}} \right)^\sigma \nonumber\\
\textit{ddfm}_{iCr}  &=& C_r \left( \textit{vdfm}_{iCr}\frac{p^C_r}{pc_{ir}} \right) \left(\frac{p^C_{ir}}{p^{\textit{dc}}_{ir}} \right)^\sigma \nonumber\\
\textit{ddfm}_{iIr}  &=& I_r \textit{vdfm}_{iIr} \nonumber\\
\textit{ddfm}_{iGr} &=& G_r \textit{vdfm}_{iGr}\left(\frac{pg_{ir}}{p^{\textit{dg}}_{ir}} \right)^\sigma   \nonumber\\
\textit{dxmd}_{isr} &=& M_{ir} \textit{vxmd}_{isr} \left(\frac{p^M_{ir}}{\textit{pyt}^m_{isr}} \right)^\sigma \nonumber\\
\textit{dst}_{jr} &=& \textit{YT}_j  \textit{vst}_{jr} \frac{p^T_j}{p^Y_{jr}}\nonumber
\end{eqnarray}

In the \SOE model ($\textit{etrndx}_{i}<\infty$) equation
(\ref{Youtput}) is replaced by two equations, one representing the
market for domestic output and another representing the market for
exports:
\begin{eqnarray}
Y_{ir} \textit{ddm}_{ir} &=&
\sum_j \textit{ddfm}_{ijr} + \textit{ddfm}_{iCr} + \textit{ddfm}_{iIr}
+ \textit{ddfm}_{iGr} \label{Doutput}\\
Y_{ir} \textit{dxm}_{ir} &=&
\sum_s \textit{dxmd}_{irs} + \textit{dst}_{ir}
\,.\label{Xoutput}
\end{eqnarray}
In these expressions, the compensated domestic and export supply
coefficients are given by:
\begin{eqnarray}
\textit{ddm}_{ir} &=& \textit{vdm}_{ir} \left( \frac{ r^Y_{ir} } {
p^D_{ir}} \right)^{\eta} \nonumber \\
\textit{dxm}_{ir} &=& \textit{vxm}_{ir} \left( \frac{ r^Y_{ir} } {
p^X_{ir}} \right)^{\eta}. \nonumber
\end{eqnarray}

\subsubsection[Private consumption]{Private
consumption (\texttt{P("c",r)})}

Consumer demand in region $r$ in the reference equilibrium is
$\textit{vom}_{Cr}$  hence:
\begin{equation}
C_r \textit{vom}_{Cr}  = \frac{\textit{RA}_r}{p^C_r}
\end{equation}

\subsubsection[Composite imports]{Composite imports (\texttt{PM(i,r)})}

The aggregate value of imports of good $i$ in region $r$ in the
reference equilibrium is $\textit{vim}_{ir}$:
\begin{equation}
M_{ir} \textit{vim}_{ir} = \sum_j \textit{difm}_{ijr} +
\textit{difm}_{iCr} + \textit{difm}_{iGr}
\end{equation}
where compensated demand functions are given by:
\begin{eqnarray}
\textit{difm}_{ijr} &=& Y_{jr} \textit{vifm}_{ijr} \left( \frac{ci_{ijr}}{p^i_{ijr}} \right)^\sigma\nonumber\\
\textit{difm}_{iCr} &=& C_r \textit{vifm}_{iCr}  \left(\frac{pc_{ir}}{p^{ic}_{ir}}\right)^\sigma  \frac{p^C_r}{\textit{pc}_{ir}} \nonumber\\
\textit{difm}_{iGr} &=& G_r \textit{vifm}_{iGr} \left(\frac{pg_{ir}}{p^{ig}_{ir}}\right)^\sigma \nonumber
\end{eqnarray}

\subsubsection[Transport services]{Transport services (\texttt{PT(j)})}

The aggregate demand (and supply) for transport service $j$ in the
benchmark equilibrium is $\textit{vtw}_j$:
\begin{equation}
\textit{YT}_j  \textit{vtw}_j = \sum_{isr} \textit{dtwr}_{jisr}
\end{equation}
where
$$\textit{dtwr}_{jisr} = M_{ir} \textit{vtwr}_{jisr}  \left(
\frac{p^M_{ir}}{\textit{pyt}^m_{isr}} \right)^\sigma\,.$$

\subsubsection[Primary factors]{Primary factors (\texttt{PF(f,r)})}

The aggregate demand (and supply) of primary factor $f$ in region $r$
is $\textit{evom}_{fr}$:
\begin{equation}
\textit{evom}_{fr} = \left\{
\begin{array}{lr}
\sum_j \textit{dfm}_{fjr} & f \in m_f \\
\textit{evom}_{fr} \textit{FT}_{fr} & f \in s_f
\end{array} \right.
\end{equation}
where the demand for primary factor is given by:
$$\textit{dfm}_{fjr}  = Y_{jr} \textit{vfm}_{fjr}
\left( \frac{c^f_{jr}}{p^{pf}_{fjr}}\right)^\sigma\,.$$

\subsubsection[Specific factors]{Specific factors (\texttt{PS(f,j,r)})}

The net value of benchmark payments to factor $f$ in sector $j$ in
region $r$ is \texttt{vfm(f,j,r)}:
\begin{equation}
\textit{vfm}_{fjr} \left( \frac{p^S_{s_{f},jr}}{p^F_{s_{f},r}} \right)^\eta =
\textit{dfm}_{fjr}\end{equation}
where the demand for primary factor is written above.

\subsection[Regional budget]{Regional budget (\texttt{RA(r)})}

Private and public incomes are given by :
\begin{equation}
\textit{RA}_r = \sum_f p^F_{fr} \textit{evom}_{fr} + p^C_n
\textit{vb}_r -  p^I_{r} \textit{vom}_{Ir} - p^G_{r}
\textit{vom}_{Gr}  + \mathcal{R}_r
\end{equation}
The base year current account deficit in region $r$ is \texttt{vb(r)},
and region $r=n$ corresponds to the ``numeraire region'' who's
consumption prices denominates international capital flows (following
conventional static trade theory, we hold the current account deficit
fixed in counterfactual analysis).
Furthermore, tax revenue in region $r$ consists of output taxes,
intermediate demand taxes, factor taxes, final demand taxes, import
tariffs and export subsidies:
\begin{equation}
\begin{array}{ll}
\mathcal{R}_r = & \sum_j \mathcal{R}^o_{jr}
                 + \sum_{ij} \left( \mathcal{R}^{fd}_{ijr} +
                 \mathcal{R}^{fi}_{ijr} \right)
                 + \sum_{fj} \mathcal{R}^f_{fjr} \\
                &+ \sum_{i}  \left( \mathcal{R}^{pd}_{ir} +
                 \mathcal{R}^{pi}_{ir} + \mathcal{R}^{gd}_{ir} +
                 \mathcal{R}^{gi}_{ir}  \right)
                 - \sum_{is} \mathcal{R}^{xs}_{irs}
                 + \sum_{is} \mathcal{R}^{ms}_{isr}
\end{array}
\end{equation}
Each of these components of tax revenue can be calculated as an
ad-valorem or proportional tax rate times a market price times the
quantity demanded or produced.

Taxes related to $Y_{jr}$ include output taxes:\footnote{Tax revenues
in the GAMS codes -- both \MCP and \MGE are represented by the
\textit{macros} indicated in square brackets.}
$$\mathcal{R}^o_{jr} =  t^o_{jr} \textit{vom}_{jr} p^Y_{jr}
Y_{jr}\, \quad [\footnotesize{\texttt{REV\_TO(g,r)}}]\,,$$
tax revenue from intermediate inputs:
$$\mathcal{R}^{fd}_{ijr} =  t^{fd}_{ijr} p^Y_{ir}
\textit{ddfm}_{ijr}\,  \quad [\footnotesize{\texttt{REV\_TFD(i,j,r)}}]\,,$$
$$\mathcal{R}^{fi}_{ijr} = t^{fi}_{ijr} p^M_{ir}
\textit{difm}_{ijr}, \quad [\footnotesize{\texttt{REV\_TFI(i,j,r)}}]\,,$$
and factor tax revenue:
$$\mathcal{R}^f_{fjr} = t^{f}_{fjr} p^F_{fr} \textit{dfm}_{fjr} \quad
[\footnotesize{\texttt{REV\_TF(f,g,r)}}]\,.$$

Taxes on household consumption of domestic and
imported goods are:
$$\mathcal{R}^{pd}_{ir} = t^{pd}_{ir} p^Y_{ir} \textit{ddfm}_{iCr},
\quad [\footnotesize{\texttt{REV\_TFD(i,"C",r)}}]\,,$$
and
$$\mathcal{R}^{pi}_{ir} = t^{pi}_{ir} p^M_{ir} \textit{difm}_{iGr}
\quad [\footnotesize{\texttt{REV\_TFI(i,"C",r)}}]\,.$$

Taxes on public demand for domestic and imported
goods are:
$$\mathcal{R}^{gd}_{ir} = t^{gd}_{ir} p^Y_{ir} \textit{ddfm}_{iGr}
\quad [\footnotesize{\texttt{REV\_TFD(i,"G",R)}}]\,,$$
and
$$\mathcal{R}^{gi}_{ir} =  t^{gi}_{ir} p^M_{ir} \textit{difm}_{iGr}
\quad [\footnotesize{\texttt{REV\_TFI(i,"G",r)}}]\,.$$

Export subsidies (paid by the government in the exporting region) are:
$$\mathcal{R}^{xs}_{irs} = t^{xs}_{irs} p^Y_{ir}
\textit{dxmd}_{irs} \quad [\footnotesize{\texttt{REV\_TXS(i,r,s)}}]\,,$$
and import tariff revenues are given by:
$$\mathcal{R}^{ms}_{isr} = t^{ms}_{isr} \left(p^Y_{is}
(1-t^{xs}_{isr}) \textit{dxmd}_{isr} + \sum_j p^T_j \textit{dtwr}_{jisr}
\right) \quad [\footnotesize{\texttt{REV\_TMS(i,s,r)}}]\,.$$

%%% Local Variables:
%%% mode: LaTeX
%%% TeX-master: "documentation"
%%% End:
